\documentclass[a4paper,11pt]{ltjsarticle}
\usepackage{luatexja}
\usepackage{luatexja-fontspec}
\usepackage{amsmath,amssymb,amsthm}
\usepackage{mathtools}
\usepackage{booktabs}
\usepackage{longtable}
\usepackage{array}
\usepackage{hyperref}
\usepackage{graphicx}
\usepackage{float}
\usepackage{geometry}
\geometry{margin=25mm}

\hypersetup{
  colorlinks=true,
  linkcolor=blue,
  urlcolor=blue,
  citecolor=blue
}

\theoremstyle{definition}
\newtheorem{theorem}{定理}[section]
\newtheorem{proposition}[theorem]{命題}
\newtheorem{lemma}[theorem]{補題}
\newtheorem{corollary}[theorem]{系}
\newtheorem{definition}[theorem]{定義}
\newtheorem{remark}[theorem]{注意}
\newtheorem{observation}[theorem]{観察}

\setlength{\parindent}{1em}
\setlength{\parskip}{0.5em}

\providecommand{\tightlist}{\setlength{\itemsep}{0pt}\setlength{\parskip}{0pt}}
\begin{document}

\section{\texorpdfstring{\(R \to 0\)
極限における測地線構造の次元的解釈}{R \textbackslash to 0 極限における測地線構造の次元的解釈}}\label{r-to-0-ux6975ux9650ux306bux304aux3051ux308bux6e2cux5730ux7ddaux69cbux9020ux306eux6b21ux5143ux7684ux89e3ux91c8}

\textbf{--- 入れ子構造と次元階層の統合的記述 ---}

\begin{center}\rule{0.5\linewidth}{0.5pt}\end{center}

\textbf{著者}: 木原 範昭 (WF System Co., Ltd.)

\textbf{日付}: 2026年4月

\textbf{DOI：}
\href{https://doi.org/10.5281/zenodo.19538106}{10.5281/zenodo.19538106}

\begin{center}\rule{0.5\linewidth}{0.5pt}\end{center}

\subsection{要旨}\label{ux8981ux65e8}

本稿では、先行論文 {[}1{]}--{[}9{]}
で確立された中心投影フレームワークに基づき、\(R \to 0\)
極限における測地線構造を、論文 {[}4{]} の入れ子構造、論文 {[}6{]}
の次元階層、および論文 {[}9{]} の Schwarzschild--de Sitter
厳密解の三者から統合的に記述する。

論文 {[}7{]}
では、親主観空間上の測地線の任意の点における子主観空間への橋渡しを整理した。論文
{[}9{]} では、論文 {[}1{]} の真空 Einstein
方程式の球対称静的厳密解として Schwarzschild--de Sitter
計量を構成し、\(R \to 0\)
極限における地平面構造の消失を記述した。本稿はこれらを受けて、\(R \to 0\)
の極限に特化し、以下の観察を記述する。すなわち、\(R \to 0\)
の極限において主観空間が呈する構造は、論文 {[}6{]}
におけるゼロ次元主観空間の記述と幾何学的に対応し、かつ論文 {[}4{]} 命題
4.1 の入れ子構造によって、この極限点から任意の曲率半径 \(R' > 0\)
を持つ新たな主観空間が構成可能であることを確認する。

\textbf{明記事項}: 本稿では何も新たに証明しない。論文 {[}4{]}
の入れ子構造、論文 {[}6{]} の次元階層、論文 {[}7{]}
の測地線の橋渡し、および論文 {[}9{]} の厳密解を組み合わせて、\(R \to 0\)
極限の幾何学的意味を整理するのみである。物理学における測地線完備性、特異点定理、ブラックホール特異点などの概念との対応は本稿の対象外である。

\begin{center}\rule{0.5\linewidth}{0.5pt}\end{center}

\subsection{1. 導入}\label{ux5c0eux5165}

先行論文 {[}4{]} は曲率半径 \(R\)
の両極限における幾何学的正則性を示し、\(R \to 0\)
極限での構造的緊張を指摘するとともに、命題 4.1
において入れ子構造の存在を証明した。論文 {[}6{]}
はゼロ次元から四次元までの主観空間の幾何学的性質を整理した。論文 {[}7{]}
は親主観空間の測地線と子主観空間の測地線の橋渡しを記述した。

本稿の目的は、これら三つの結果を統合し、\(R \to 0\)
の極限における測地線構造に焦点を当てた記述を行うことにある。具体的には、以下の問いに答える。

\textbf{問い}: 親主観空間 \(\Pi_R\) 上の測地線を \(R \to 0\)
の方向に追跡したとき、極限で現れる構造は、論文 {[}6{]}
の次元階層のどの段階に対応し、そこからどのような幾何学的自由度が残存するか。

本稿は何も新たに証明しない。既存の命題の組み合わせとして記述するのみである。

\begin{center}\rule{0.5\linewidth}{0.5pt}\end{center}

\subsection{2. 準備}\label{ux6e96ux5099}

\subsubsection{2.1
用語と先行結果の要約}\label{ux7528ux8a9eux3068ux5148ux884cux7d50ux679cux306eux8981ux7d04}

本稿で使用する先行論文の結果を再掲する。

\textbf{論文 {[}4{]} 命題 4.1（入れ子構造の存在）}: 任意の主観空間
\(\Pi_R\) の内部の任意の点 \(p\) に対し、\(p\)
を中心とする新たな中心投影による主観空間 \(\Pi_{R'}\) が、任意の
\(R' > 0\) に対して存在する。\(\Pi_{R'}\) は論文 {[}1{]}
の構成と同じ幾何学的構造を持つ。

\textbf{論文 {[}4{]} §3.2（構造的緊張）}: \(R \to 0\)
の極限において、測地線偏差の上限 \(|\xi|^2_{\max}(R)\)
はゼロに漸近する一方、局所曲率は \(1/R^2\) のオーダーで発散する。

\textbf{論文 {[}6{]} §2（ゼロ次元主観空間）}: ゼロ次元の主観空間は「ある
/
ない」のみを区別する世界であり、内部構造を持たず、測地線偏差は存在せず、自身が曲率半径に支配されているかどうかを内部から認知する手段も存在しない。

\textbf{論文 {[}7{]} §4（測地線の橋渡し）}: 親主観空間 \(\Pi_R\)
上の測地線 \(\gamma\) の任意の点 \(p\) を立ち上がり点として子主観空間
\(\Pi_{R'}\) が構成可能であり、\(\Pi_{R'}\) 上でも測地線が独立に
well-defined である。

\subsubsection{2.2
本稿の位置付け}\label{ux672cux7a3fux306eux4f4dux7f6eux4ed8ux3051}

論文 {[}7{]}
は、親主観空間上の測地線の\textbf{任意の点}における子主観空間への橋渡しを一般的に記述した。本稿は、この一般的記述を
\(R \to 0\) の\textbf{極限}に特化して適用し、論文 {[}6{]}
の次元階層との対応を加えることで、極限の幾何学的意味を整理する。

\begin{center}\rule{0.5\linewidth}{0.5pt}\end{center}

\subsection{\texorpdfstring{3. \(R \to 0\)
極限とゼロ次元主観空間の対応}{3. R \textbackslash to 0 極限とゼロ次元主観空間の対応}}\label{r-to-0-ux6975ux9650ux3068ux30bcux30edux6b21ux5143ux4e3bux89b3ux7a7aux9593ux306eux5bfeux5fdc}

\subsubsection{3.1
測地線偏差の消失と次元的解釈}\label{ux6e2cux5730ux7ddaux504fux5deeux306eux6d88ux5931ux3068ux6b21ux5143ux7684ux89e3ux91c8}

論文 {[}4{]} §3.2 により、\(R \to 0\) の極限では測地線偏差の上限
\(|\xi|^2_{\max}(R) \to 0\)
である。すなわち、この極限において内部観測者が測定可能な偏差はゼロに漸近する。

一方、論文 {[}6{]} §2
のゼロ次元主観空間の記述では、測地線偏差が「存在しない」とされている。

この二つの記述は、異なる経路から同一の構造的特徴に至る。前者は \(n\)
次元主観空間の \(R \to 0\)
極限として、後者はゼロ次元主観空間の内在的性質として、いずれも「測地線偏差が観測不能である」という状態を記述している。

\subsubsection{3.2
内部構造の消失と次元的解釈}\label{ux5185ux90e8ux69cbux9020ux306eux6d88ux5931ux3068ux6b21ux5143ux7684ux89e3ux91c8}

論文 {[}6{]} §2
では、ゼロ次元主観空間の内部構造は「存在しない」とされ、自身が曲率半径
\(R\)
に支配されているかどうかを内部から認知する手段も存在しないとされた。

\(R \to 0\) の極限において、主観空間 \(\Pi_R\) は「点に縮む」（論文
{[}4{]} 問い
3.2）。点に縮んだ主観空間の内部観測者は、方向の区別も距離の測定もできない。これは論文
{[}6{]} §2 のゼロ次元主観空間の記述と幾何学的に一致する。

\subsubsection{3.3 対応の記述}\label{ux5bfeux5fdcux306eux8a18ux8ff0}

以上をまとめると、次の観察が得られる。

\textbf{観察 3.1}: \(n\) 次元主観空間 \(\Pi_R\) の \(R \to 0\)
極限は、論文 {[}6{]}
におけるゼロ次元主観空間の記述と以下の点で幾何学的に対応する。

\begin{enumerate}
\def\labelenumi{\arabic{enumi}.}
\tightlist
\item
  測地線偏差の測定可能な値がゼロに漸近する（論文 {[}4{]}
  §3.2）ことは、ゼロ次元主観空間において測地線偏差が存在しない（論文
  {[}6{]} §2）ことに対応する。
\item
  主観空間が点に縮む（論文 {[}4{]} 問い
  3.2）ことは、ゼロ次元主観空間が内部構造を持たない（論文 {[}6{]}
  §2）ことに対応する。
\item
  内部観測者が曲率半径を認知できなくなることは、ゼロ次元主観空間において曲率半径の認知手段が存在しない（論文
  {[}6{]} §2）ことに対応する。
\end{enumerate}

この観察は新たな証明を含まない。論文 {[}4{]} と論文 {[}6{]}
の既存の記述を並置したものである。

\begin{center}\rule{0.5\linewidth}{0.5pt}\end{center}

\subsection{4.
極限点からの子主観空間の立ち上げ}\label{ux6975ux9650ux70b9ux304bux3089ux306eux5b50ux4e3bux89b3ux7a7aux9593ux306eux7acbux3061ux4e0aux3052}

\subsubsection{4.1
入れ子構造の適用}\label{ux5165ux308cux5b50ux69cbux9020ux306eux9069ux7528}

論文 {[}4{]} 命題 4.1 により、主観空間 \(\Pi_R\) の内部の任意の点 \(p\)
から子主観空間 \(\Pi_{R'}\) を立ち上げることができる。\(R > 0\)
である限り、\(p\) は \(\Pi_R\) の内部点であり、命題 4.1
の前提条件を満たす。

\(R \to 0\) の極限において、\(\Pi_R\) は §3
で述べたようにゼロ次元的な構造に漸近する。しかし、\(R > 0\)
の任意の値に対して命題 4.1 は成立し、立ち上げられる子主観空間
\(\Pi_{R'}\) の曲率半径 \(R'\) は \(R\)
とは独立に任意の正の値をとりうる。

\subsubsection{4.2
次元階層との接続}\label{ux6b21ux5143ux968eux5c64ux3068ux306eux63a5ux7d9a}

§3 の観察 3.1 と命題 4.1 を組み合わせると、次の記述が得られる。

\textbf{観察 4.1}: 親主観空間 \(\Pi_R\) が \(R \to 0\)
でゼロ次元的な構造に漸近する過程において、その漸近領域の任意の点から、論文
{[}4{]} 命題 4.1 により、任意の \(R' > 0\) を持つ子主観空間 \(\Pi_{R'}\)
が構成可能である。子主観空間 \(\Pi_{R'}\) は論文 {[}1{]}
の構成と同じ幾何学的構造を持ち、その内部では論文 {[}6{]}
の記述に従い、\(R'\) の値に応じた次元の幾何学的性質が展開される。

この観察を論文 {[}6{]} の次元階層の文脈で言い換えると:

\begin{itemize}
\tightlist
\item
  親主観空間のゼロ次元的極限は、論文 {[}6{]} §2 の「ある /
  ない」のみの世界に対応する。
\item
  しかし、その「ある」の側から、論文 {[}4{]} 命題 4.1
  により、新たな主観空間が任意の \(R' > 0\)
  で立ち上がる幾何学的自由度が存在する。
\item
  立ち上がった子主観空間 \(\Pi_{R'}\) は、\(R'\) の値に応じて論文
  {[}6{]} §3 以降の次元的構造を持つ。
\end{itemize}

\subsubsection{4.3
本節の主張の範囲}\label{ux672cux7bc0ux306eux4e3bux5f35ux306eux7bc4ux56f2}

本節で述べたのは以下に限定される。

\begin{enumerate}
\def\labelenumi{\arabic{enumi}.}
\tightlist
\item
  \(R \to 0\) の極限領域においても、\(R > 0\) である限り命題 4.1
  が適用可能であること。
\item
  立ち上げられた子主観空間の曲率半径 \(R'\) は親の \(R\)
  とは独立であること。
\item
  これらの組み合わせにより、親のゼロ次元的極限と子の有限次元的構造が共存する記述が可能であること。
\end{enumerate}

本節は新たな証明を含まない。命題 4.1 と論文 {[}6{]}
の記述の組み合わせである。

\begin{center}\rule{0.5\linewidth}{0.5pt}\end{center}

\subsection{5.
測地線の追跡と次元的転換}\label{ux6e2cux5730ux7ddaux306eux8ffdux8de1ux3068ux6b21ux5143ux7684ux8ee2ux63db}

\subsubsection{\texorpdfstring{5.1 親主観空間における測地線の
\(R \to 0\)
方向への追跡}{5.1 親主観空間における測地線の R \textbackslash to 0 方向への追跡}}\label{ux89aaux4e3bux89b3ux7a7aux9593ux306bux304aux3051ux308bux6e2cux5730ux7ddaux306e-r-to-0-ux65b9ux5411ux3078ux306eux8ffdux8de1}

論文 {[}7{]} §3 で記述されたように、親主観空間 \(\Pi_R\) 上の測地線
\(\gamma\) を追跡すると、\(R \to 0\) の方向において測地線偏差の上限
\(|\xi|^2_{\max}(R)\) がゼロに漸近する。

§3 の観察 3.1
により、この漸近は親主観空間のゼロ次元的縮退として次元的に解釈できる。すなわち、親主観空間の内部観測者にとって、\(\gamma\)
の追跡は次元的構造が失われる過程として現れる。

\subsubsection{5.2
子主観空間における測地線の再出発}\label{ux5b50ux4e3bux89b3ux7a7aux9593ux306bux304aux3051ux308bux6e2cux5730ux7ddaux306eux518dux51faux767a}

論文 {[}7{]} §4 により、\(\gamma\) 上の点 \(p\)
を立ち上がり点とする子主観空間 \(\Pi_{R'}\) において、\(p\)
から出発する測地線 \(\gamma'\) が well-defined である。§4 の観察 4.1
により、\(R \to 0\)
の極限領域においてもこの立ち上げは可能であり、\(\Pi_{R'}\) 上の
\(\gamma'\) は \(R'\) に応じた次元的構造の中で定義される。

\subsubsection{5.3
次元的転換としての記述}\label{ux6b21ux5143ux7684ux8ee2ux63dbux3068ux3057ux3066ux306eux8a18ux8ff0}

§5.1 と §5.2 を組み合わせると、次の記述が得られる。

\textbf{観察 5.1}: 親主観空間 \(\Pi_R\) 上の測地線 \(\gamma\) を
\(R \to 0\)
の方向に追跡する過程は、以下の二つの段階として幾何学的に記述できる。

\textbf{段階 1（次元的縮退）}: 親主観空間 \(\Pi_R\)
において、測地線偏差の上限 \(|\xi|^2_{\max}(R)\)
がゼロに漸近し、主観空間は論文 {[}6{]} §2
のゼロ次元的構造に近づく。親の内部観測者にとって、測地線の追跡は次元的構造の喪失として現れる。

\textbf{段階 2（次元的再構成）}: 論文 {[}4{]} 命題 4.1
により、極限領域の点から子主観空間 \(\Pi_{R'}\) が任意の \(R' > 0\)
で構成可能であり、子主観空間の内部では論文 {[}6{]}
の次元階層に従った幾何学的構造が展開される。子の内部観測者にとって、測地線
\(\gamma'\) は \(R'\) に応じた次元的構造の中で well-defined である。

この二段階の記述において、親主観空間における測地線の「終端」は、幾何学的構造の消滅ではなく、次元的構造の転換点として記述される。親の内部観測者から見た構造の消失と、子の内部観測者から見た構造の存在は、論文
{[}3{]} 定理 5.1（観測の相対性）によって矛盾なく共存する。

\subsubsection{5.4
本節の主張の範囲}\label{ux672cux7bc0ux306eux4e3bux5f35ux306eux7bc4ux56f2-1}

本節で述べたのは以下に限定される。

\begin{enumerate}
\def\labelenumi{\arabic{enumi}.}
\tightlist
\item
  §3 の次元的対応と §4 の入れ子構造の組み合わせにより、\(R \to 0\)
  極限における測地線の振る舞いが「次元的転換」として記述可能であること。
\item
  この記述は親と子の観測者の双方にとって整合的であること（論文 {[}3{]}
  定理 5.1 による）。
\item
  「次元的転換」は幾何学的記述のための用語であり、物理的過程を主張するものではないこと。
\end{enumerate}

本節は新たな証明を含まない。

\begin{center}\rule{0.5\linewidth}{0.5pt}\end{center}

\subsection{6. 考察}\label{ux8003ux5bdf}

本稿では、\(R \to 0\) 極限における測地線構造を、論文 {[}4{]}
の入れ子構造、論文 {[}6{]} の次元階層、論文 {[}7{]}
の測地線の橋渡し、および論文 {[}9{]} の Schwarzschild--de Sitter
厳密解の四者を統合する形で記述した。

本稿の記述は以下に限定される。

\begin{enumerate}
\def\labelenumi{\arabic{enumi}.}
\tightlist
\item
  \(R \to 0\) 極限が論文 {[}6{]}
  のゼロ次元主観空間の記述と幾何学的に対応すること（観察 3.1）。
\item
  ゼロ次元的極限の領域から子主観空間が構成可能であり、子の内部では次元的構造が展開されること（観察
  4.1）。
\item
  親主観空間の測地線の \(R \to 0\)
  方向への追跡が、次元的縮退と次元的再構成の二段階として記述可能であること（観察
  5.1）。
\end{enumerate}

論文 {[}9{]}
は、この次元的転換の過程を厳密解のレベルで裏付ける。すなわち、\(R \to 0\)
で宇宙論的地平面が消失し内部観測者の領域が失われる過程（論文 {[}9{]}
§4）が、本稿の「段階1（次元的縮退）」に対応し、子主観空間における地平面構造の回復（論文
{[}9{]} §5）が「段階2（次元的再構成）」に対応する。

さらに、段階2で立ち上がる子主観空間の安定次元について、論文 {[}8{]}
が組合せ論的な考察を提供している。論文 {[}8{]}
は、ゼロ次元の構造的情報から出発して、離散性と安定性の仮定のもとで四次元が最小の安定次元として選ばれることを示した。本稿の観察
5.1 における「次元的再構成」の到達先が論文 {[}8{]} の議論と接続する。

本稿で記述した「次元的転換」は、幾何学の枠組みの内部における記述法であり、物理的過程の主張ではない。特に、定説とされる微分幾何学上の測地線完備性の概念、一般相対論における特異点定理、ブラックホール時空における測地線の振る舞いなどとの対応は、幾何学の枠組みの外側の問題であり、本稿の射程を超える。

本稿の記述が持ちうる物理的意味付けについては、読者の判断に委ねる。

\begin{center}\rule{0.5\linewidth}{0.5pt}\end{center}

\subsection{7. 結論}\label{ux7d50ux8ad6}

中心投影フレームワークにおいて、主観空間 \(\Pi_R\) の \(R \to 0\)
極限は、論文 {[}6{]} のゼロ次元主観空間の記述と幾何学的に対応する。論文
{[}4{]} 命題 4.1 の入れ子構造により、このゼロ次元的極限の領域から任意の
\(R' > 0\)
を持つ子主観空間が構成可能であり、子の内部では次元的構造が展開される。親主観空間上の測地線の
\(R \to 0\)
方向への追跡は、次元的縮退と次元的再構成の二段階として記述され、測地線の「終端」は幾何学的構造の消滅ではなく次元的構造の転換点として整理される。これらはすべて論文
{[}1{]}--{[}7{]}
の既存の命題の組み合わせとして記述される。物理的解釈は本稿の対象外であり、読者に委ねる。

\begin{center}\rule{0.5\linewidth}{0.5pt}\end{center}

\subsection{参考文献}\label{ux53c2ux8003ux6587ux732e}

{[}1{]} 木原 範昭,「中心投影による4次元空間の幾何学的定式化」, Zenodo,
DOI: 10.5281/zenodo.19427780 (2026).

{[}2{]} 木原 範昭,「中心投影の幾何学的対称性：多軸モデルの数学的基盤」,
Zenodo, DOI: 10.5281/zenodo.19434932 (2026).

{[}3{]} 木原
範昭,「複数の主観空間における観測の相対性：中心投影の対称性の幾何学的帰結」,
Zenodo, DOI: 10.5281/zenodo.19435162 (2026).

{[}4{]} 木原 範昭,「中心投影における曲率半径の極限についての考察 ---
\(R \to \infty\) と \(R \to 0\) の場合」, Zenodo, DOI:
10.5281/zenodo.19526549 (2026).

{[}5{]} 木原 範昭,「背景空間と主観空間への次元追加についての考察」,
Zenodo, DOI: 10.5281/zenodo.19526913 (2026).

{[}6{]} 木原 範昭,「ゼロ次元から四次元主観空間の考察」, Zenodo, DOI:
10.5281/zenodo.19533292 (2026).

{[}7{]} 木原 範昭,「主観空間における測地線の連続性についての考察 ---
親主観空間から子主観空間への測地線の橋渡しについて」, Zenodo, DOI:
10.5281/zenodo.19533299 (2026).

{[}8{]} 木原 範昭,「体積1の四次元超直方体に外接する超球体の直径」,
Zenodo, DOI: 10.5281/zenodo.19533313 (2026).

{[}9{]} 木原 範昭,「中心投影フレームワークにおける Schwarzschild--de
Sitter 厳密解 --- 論文 {[}1{]} の真空 Einstein
方程式の球対称静的厳密解」, Zenodo, DOI: 10.5281/zenodo.19538098 (2026).

\begin{center}\rule{0.5\linewidth}{0.5pt}\end{center}

\end{document}
